\section{結果}
サリチル酸 \SI{2.00}{g} を用いて反応を行い，再結晶後にアセチルサリチル酸 \SI{2.17}{g} を得た．
サリチル酸のモル質量を \SI{138.12}{g/mol}，アセチルサリチル酸のモル質量を \SI{180.16}{g/mol} とすると，サリチル酸の物質量は $2.00 / 138.12 = \SI{1.45e-2}{mol}$ である．
したがって理論収量は $\SI{1.45e-2}{mol} \times \SI{180.16}{g/mol} = \SI{2.61}{g}$ となり，収率は次式で計算できる．

\begin{equation}
\text{収率} = \frac{\SI{2.17}{g}}{\SI{2.61}{g}} \times 100 = 83.1\,
\%
\label{eq:percent-yield}
\end{equation}

\begin{table}[H]
    \centering
    \caption{アセチルサリチル酸合成の結果}
    \label{tab:aspirin-result}
    \begin{tabular}{lcc}
        \hline
        項目 & 値 & 備考 \\
        \hline
        サリチル酸質量 & \SI{2.00}{g} & 制限試薬 \\
        再結晶後質量 & \SI{2.17}{g} & 白色結晶 \\
        理論収量 & \SI{2.61}{g} & 1:1 量論比 \\
        収率 & 83.1\% & 式\eqnref{eq:percent-yield} \\
        融点 & \SIrange{134.5}{135.8}{\celsius} & 文献値 \SIrange{135}{136}{\celsius} \\
        \hline
    \end{tabular}
\end{table}

\figref{fig:aspirin-ring} に示すように，生成物ではベンゼン環上のヒドロキシ基がアセチル化されている．得られた結晶は白色針状であり，文献値に近い融点を示したことから，目的物が比較的高純度で得られたと考えられる．